\chapter{Sugestões, Comentários e Propostas de Melhoria}\label{chap5}
\thispagestyle{fancy}

Neste capítulo é realizada uma análise reflexiva e crítica sobre o decorrer do estágio curricular, integrando as sugestões organizacionais, comentários sobre a experiência de trabalho, as dificuldades enfrentadas, as limitações identificadas no sistema e as propostas de trabalho futuro.

\section{Sugestões}

Relativamente ao funcionamento do estágio e à sua articulação académica, sugere-se maior flexibilidade na definição inicial de cronogramas para projetos com integração de múltiplos módulos. Projetos que envolvem quiosques para ecrãs táteis e controlo de fluxos logísticos estão sujeitos a adaptações frequentes e reestruturação de bases de dados locais e remotas. Um período de margem oficial para testes de integração em ambiente real diminuiria riscos no planeamento.

\section{Comentários}

A integração na 20 Recolher revelou-se uma experiência importante de aprendizagem profissional. O contacto direto com o setor de gestão de resíduos permitiu compreender a importância da digitalização e da automatização de processos logísticos críticos. A colaboração estreita entre os dois estagiários facilitou a discussão de decisões de arquitetura e fomentou boas práticas de trabalho em equipa, controlo de versões e documentação de APIs.

\section{Dificuldades Encontradas e Soluções Adotadas}

Durante o ciclo de desenvolvimento, a equipa enfrentou três dificuldades principais:
\begin{itemize}
    \item \textbf{Gestão de Estados Concorrentes nos Pedidos:} A existência de múltiplos acessos simultâneos por parte de clientes e operadores gerava inconsistências temporárias no estado das recolhas (por exemplo, dois operadores a tentarem aprovar o mesmo pedido). Para resolver isto, foi implementada uma validação ao nível da API baseada em controlo de concorrência otimista e tratamento de exceções no EF Core.
    \item \textbf{Migração de Arquitetura Monolítica para Distribuída:} Numa fase inicial, o backoffice e o terminal de pesagem foram idealizados numa única aplicação Blazor. Contudo, isso gerava duplicação de acessos à base de dados. Adotou-se como solução isolar o acesso aos dados numa API central baseada em \textit{controllers} RESTful. Esta mudança exigiu a refatoração da comunicação para chamadas HTTP assíncronas em JSON, o que estabilizou o sistema.
    \item \textbf{Gestão de Cache no Supabase e Atualizações Realtime:} Garantir que as notificações de chat e o estado dos pedidos fossem atualizados sem sobrecarregar a base de dados com consultas repetitivas (\textit{polling}) foi complexo. A solução consistiu em configurar o cliente Realtime do Supabase baseado em WebSockets apenas para eventos específicos (mensagens de chat), mantendo o restante fluxo sob pedidos REST padrão.
\end{itemize}

\section{Limitações do Projeto}

Apesar de o ecossistema tecnológico ter atingido os objetivos estipulados, persistem algumas limitações:
\begin{itemize}
    \item \textbf{Dependência de Alojamento Cloud (Supabase/Vercel):} Dado que a base de dados centralizada e o website estão alojados em infraestruturas cloud externas, o ecossistema necessita de ligação permanente à Internet. No caso de uma falha de rede física no armazém da empresa, o quiosque de pesagem não conseguirá registar novas pesagens de forma automática na API.
    \item \textbf{Falta de Módulo de Faturação Direta Integrado:} Atualmente, o sistema extrai listagens em formatos estruturados (CSV/PDF) para permitir a faturação manual. O ecossistema não possui uma integração direta via Web Services com o software de faturação certificado da empresa (ex: Primavera ou Sage).
\end{itemize}

\section{Propostas de Melhoria (Trabalho Futuro)}

Para futuros desenvolvimentos e evolução contínua da plataforma, propõem-se as seguintes melhorias técnicas:
\begin{itemize}
    \item \textbf{Modo Offline com Sincronização Posterior no Quiosque:} Implementar um mecanismo de persistência de pesagens local no terminal (usando bases de dados SQLite locais ou armazenamento local encriptado) que permita continuar a operar sem ligação à Internet e envie os dados acumulados de forma assíncrona assim que a ligação à API for restabelecida.
    \item \textbf{Integração de APIs de Software de Faturação Certificado:} Desenvolver conectores dedicados para comunicar com o ERP da empresa, automatizando a emissão imediata da guia de transporte e da fatura após a validação da pesagem no terminal quiosque.
    \item \textbf{Notificações Push e Aplicação Móvel para Clientes:} Expandir a presença web do cliente para uma aplicação híbrida móvel ou implementar notificações push nativas no navegador, alertando instantaneamente os utilizadores quando uma recolha for efetuada ou quando for recebida resposta do suporte técnico.
\end{itemize}
